\documentclass{article}

\usepackage{fontawesome5}
\usepackage{graphicx}
\usepackage{hyperref}
\usepackage{multicol}
\usepackage{geometry}
 \geometry{
 a4paper,
 total={170mm,257mm},
 left=20mm,
 top=40mm,
 }
 
 
\begin{document}

\section*{\Huge Ravi Chandra Reddy N}

\begin{tabular}{lcr}
\begin{tabular}{lc}
\faIcon{linkedin} \href{https://linkedin.com/in/rcreddyn}{LinkedIn} & \faIcon{globe} \href{https://rcreddyn.github.io}{https://rcreddyn.github.io}\\
\end{tabular}

\begin{tabular}{cr}
\faIcon{phone} +919490977952 & \faIcon{envelope} rcreddy1997@gmail.com   \\
\end{tabular}
\end{tabular}

\begin{multicols}{2}
\section*{Education}
\faIcon{university} BV Raju Institute of Technology (2014 - 2018)
\begin{itemize}
\item \small Bachelor of Technology in Electrical and Electronics Engineering with 76\%.
\item \small Organizer for Technical, Cultural events and Industrial visits.
\item \small Supplementary education in Embedded and Assistive Technology  at Assitive Technology Lab.
\end{itemize}

\section*{Skills}
\begin{itemize}
    \item Java, C, and Basic C++
    \item  Python, and Shell Scripting
    \item Algorithms, and Data Structures
    \item PostgreSQL
    \item HTML, CSS
    \item Git, and \LaTeX
\end{itemize}

\section*{Coding profile(s)}
\begin{itemize}
\item \href {https://hackerrank.com/rcreddyn}{HackerRank}
\item \href {https://auth.geeksforgeeks.org/user/rcreddyn/}{GeeksforGeeks}
\end{itemize}

\section*{Experience}
\includegraphics[width=14px]{isro.png} National Remote Sensing Center (2018)
\begin{itemize}
\item \small Project Student
\item \small Designed and developed standard ways to trans-receive data from sensors and actuators to  Bhuvan IoT Cloud.
\end{itemize}

\section*{Projects}
\begin{itemize}
\item \href{https://github.com/rcreddyn/nomsh}{Nomsh}, a shell written in C.
\begin{itemize}
\item Supports execution of builtin, and executable commands.
\item Supports output redirection.
\item Wrote a Makefile for creating an executable.
\end{itemize}
\item  \href{https://github.com/rcreddyn/nomsh}{Amrika}, a compiler written in C++ for a cooked up language.
\begin{itemize}
\item Wrote grammar for a language, and implemented parser, lexer, and emiter from scratch.
\item Implemented print and assignment statements and analogues to if conditional and while iteration.
\item Wrote a makefile to generate the compiler, and a bash script to run the compiled code.
\end{itemize}
\item \href{https://github.com/rcreddyn/nomsh}{Lavangam}, a trie-based command line spell checker written in Python.
\begin{itemize}
\item Implemented a trie to store/search for words.
\item Implemented Levenshtein's algorithm to calculate distance between two words.
\item Utilized wordlist from NLTK corpus.
\end{itemize}
\item \href{https://github.com/rcreddyn/nomsh}{Holyperil}, a command line planetory monitering application written in Python. 
\begin{itemize}
\item Code written in Python to notify of potentially hazardous/ dangerous asteroids approaching Earth.
\item Used an offset of six days, to prepare just in case of an approaching hypothetical danger.
\item Utilized NASA's Asteroids - NeoWs API for the search.
\end{itemize}
\item \href{https://github.com/rcreddyn/nomsh}{Freddie}, a text-only browser.
\begin{itemize}
\item Code written in Python to extract text from webpages, bypassing HTML tags, CSS, and all kinds of media.
\end{itemize}
\end{itemize} 
\end{multicols}
\end{document}